% Options for packages loaded elsewhere
\PassOptionsToPackage{unicode}{hyperref}
\PassOptionsToPackage{hyphens}{url}
\documentclass[
  11pt]{article}
\usepackage{xcolor}
\usepackage[margin=1in]{geometry}
\usepackage{amsmath,amssymb}
\setcounter{secnumdepth}{-\maxdimen} % remove section numbering
\usepackage{iftex}
\ifPDFTeX
  \usepackage[T1]{fontenc}
  \usepackage[utf8]{inputenc}
  \usepackage{textcomp} % provide euro and other symbols
\else % if luatex or xetex
  \usepackage{unicode-math} % this also loads fontspec
  \defaultfontfeatures{Scale=MatchLowercase}
  \defaultfontfeatures[\rmfamily]{Ligatures=TeX,Scale=1}
\fi
\usepackage{lmodern}
\ifPDFTeX\else
  % xetex/luatex font selection
\fi
% Use upquote if available, for straight quotes in verbatim environments
\IfFileExists{upquote.sty}{\usepackage{upquote}}{}
\IfFileExists{microtype.sty}{% use microtype if available
  \usepackage[]{microtype}
  \UseMicrotypeSet[protrusion]{basicmath} % disable protrusion for tt fonts
}{}
\makeatletter
\@ifundefined{KOMAClassName}{% if non-KOMA class
  \IfFileExists{parskip.sty}{%
    \usepackage{parskip}
  }{% else
    \setlength{\parindent}{0pt}
    \setlength{\parskip}{6pt plus 2pt minus 1pt}}
}{% if KOMA class
  \KOMAoptions{parskip=half}}
\makeatother
% definitions for citeproc citations
\NewDocumentCommand\citeproctext{}{}
\NewDocumentCommand\citeproc{mm}{%
  \begingroup\def\citeproctext{#2}\cite{#1}\endgroup}
\makeatletter
 % allow citations to break across lines
 \let\@cite@ofmt\@firstofone
 % avoid brackets around text for \cite:
 \def\@biblabel#1{}
 \def\@cite#1#2{{#1\if@tempswa , #2\fi}}
\makeatother
\newlength{\cslhangindent}
\setlength{\cslhangindent}{1.5em}
\newlength{\csllabelwidth}
\setlength{\csllabelwidth}{3em}
\newenvironment{CSLReferences}[2] % #1 hanging-indent, #2 entry-spacing
 {\begin{list}{}{%
  \setlength{\itemindent}{0pt}
  \setlength{\leftmargin}{0pt}
  \setlength{\parsep}{0pt}
  % turn on hanging indent if param 1 is 1
  \ifodd #1
   \setlength{\leftmargin}{\cslhangindent}
   \setlength{\itemindent}{-1\cslhangindent}
  \fi
  % set entry spacing
  \setlength{\itemsep}{#2\baselineskip}}}
 {\end{list}}
\usepackage{calc}
\newcommand{\CSLBlock}[1]{\hfill\break\parbox[t]{\linewidth}{\strut\ignorespaces#1\strut}}
\newcommand{\CSLLeftMargin}[1]{\parbox[t]{\csllabelwidth}{\strut#1\strut}}
\newcommand{\CSLRightInline}[1]{\parbox[t]{\linewidth - \csllabelwidth}{\strut#1\strut}}
\newcommand{\CSLIndent}[1]{\hspace{\cslhangindent}#1}
\setlength{\emergencystretch}{3em} % prevent overfull lines
\providecommand{\tightlist}{%
  \setlength{\itemsep}{0pt}\setlength{\parskip}{0pt}}
\usepackage{amsmath,amssymb,amsthm}
\usepackage{booktabs}
\usepackage{graphicx}
\usepackage{float}
\usepackage{hyperref}
\allowdisplaybreaks
\usepackage{bookmark}
\IfFileExists{xurl.sty}{\usepackage{xurl}}{} % add URL line breaks if available
\urlstyle{same}
\hypersetup{
  pdftitle={Competition Creates Competition: Stock Prices and the Discovery of Takeover Bidders},
  hidelinks,
  pdfcreator={LaTeX via pandoc}}

\title{Competition Creates Competition: Stock Prices and the Discovery
of Takeover Bidders}
\author{}
\date{}

\begin{document}
\maketitle

\textbf{Austin Li}

\section{Abstract}\label{abstract}

A stronger acquirer can attract its own challenger. We study a listed
target whose prospective buyer learns from stock prices before paying
for acquisition diligence. Strengthening an incumbent reduces the
challenger's acquisition profit at every fixed belief, but increases the
sensitivity of target shareholders' proceeds to challenger quality. This
induces informed trading and can reverse entry deterrence. Once the
price-information experiment is fixed, conventional deterrence
re-emerges. We characterize unique uninformative and informative
outcomes and establish coexistence with asymmetric trading. The
mechanism survives smooth noise, atomless preparation costs, and
complementary signals: the buyer can have the more accurate signal.
Holding competitive strength fixed, access to prices raises acquisition
surplus net of diligence. A bargaining comparison identifies the payment
property behind discovery. We formulate the seller's design problem when
sale terms determine both acquisition payments and the information that
brings buyers into the contest.

\section{1. Introduction}\label{sec-introduction}

Where does takeover competition come from? A prospective acquirer must
decide whether an opportunity warrants investigation before it knows
enough to submit an executable bid. A powerful rival makes that
investigation less attractive: it reduces the probability of a
profitable acquisition and the rents available if the challenger wins.
Yet the same rival can make the target's shares more informative about
whether the challenger should investigate. We show that this second
force can dominate the first. \textbf{A stronger acquirer can recruit
the challenger it would ordinarily deter.}

The distinction is between the return to acquiring a company and the
return to trading its shares. When an excellent challenger faces weak
competition, much of its acquisition advantage remains with the
challenger. Its quality matters for the value created by a transaction
but need not substantially change what target shareholders receive.
Stronger competition transfers more of that advantage into the target's
sale price. Information about the challenger consequently becomes more
valuable to investors trading target equity. Their orders reveal
information that the challenger uses before committing resources to
diligence and participation.

This mechanism joins a corporate-control contest to a financial market
in which prices are rational and traders anticipate the real decisions
their orders influence. It does not rely on an unanticipated takeover
premium. The market maker prices the entry response. What remains
profitable to trade on is the difference between high- and low-quality
acquisition outcomes after that response has been priced.

We begin with a committed cash auction whose acquisition payments are
disciplined by the strongest competing bid. An incumbent is already
prepared to participate. A potential challenger learns from the target's
price before paying a privately realized preparation cost and evaluating
its acquisition value. An informed investor trades the target's shares
against noise demand. The incumbent's value distribution is the
comparative-static primitive. We derive the acquisition payoffs, solve
inference from the price actually observed by the challenger, and verify
global trading incentives over continuous orders, allowing mixed
strategies.

The main result compares regions with different equilibrium information.
With a weak incumbent, the information-sensitive spread in target
proceeds is too small to support informed trading. With a stronger
incumbent, informed trading is uniquely optimal and sufficiently
favorable prices induce additional challenger entry. Acquisition profits
decrease at every fixed posterior throughout this comparison. Entry
rises because the information supplied by the market changes.

The equilibrium correspondence contains additional economics. When
orders are fixed, the price-information experiment is fixed, and
stronger competition weakly reduces entry. Before orders reach their
bounds, asymmetric trading can change the experiment: buying after
favorable information and selling after unfavorable information need not
have equal magnitudes. We establish informative asymmetric equilibria
with increasing entry across certified intermediate strengths,
coexisting with an uninformative equilibrium. Our correspondence display
preserves these distinct outcomes rather than selecting a curve. At
still higher strengths, the bounded likelihood ratios of the benchmark
noise law can make expensive entry infeasible even though trading
remains informative.

The information interpretation is complementary knowledge, not a
presumption that outsiders know more than prospective buyers. We give
the buyer and the investor different noisy signals. The buyer can have
the more accurate signal and still use the incremental information in
prices. This extension rederives conditional participation and
competitive pricing when the buyer's private information predicts entry
even after the public price is observed.

We make the following contributions.

\begin{enumerate}
\def\labelenumi{\arabic{enumi}.}
\tightlist
\item
  \textbf{Competition reallocates the returns to information.} We derive
  an acquisition-stage opposition, general in the incumbent's value
  distribution: stronger competition raises the information sensitivity
  of target proceeds while lowering the challenger's acquisition profit
  at every posterior.
\item
  \textbf{Prices can reverse entry deterrence.} We prove the equilibrium
  entry reversal, the sufficiency of the observed price, and global best
  responses. We distinguish a changing information experiment from a
  fixed informative experiment and characterize both uniquely determined
  and coexisting outcomes.
\item
  \textbf{Discovery has institutional and real consequences.} We
  establish complementary-signal and smooth-distribution extensions, a
  matched surplus gain from access to prices, and a bargaining boundary
  identifying when runner-up-driven pricing supports the payoff
  opposition.
\item
  \textbf{Sale design becomes a problem of attracting information as
  well as buyers.} We specify the seller's continuation game and
  objective over the full reserve domain. The next result is a
  characterization of seller-optimal discovery and the role of
  commitment before information is revealed.
\end{enumerate}

The acquisition mechanism is fixed in the entry theorem; the seller's
choice is a separate optimization problem. That separation permits a
sharp first result and identifies the next economic question. A seller
choosing sale terms changes not only what an entrant pays but also
whether investors have a reason to reveal the information that makes
entry worthwhile.

\section{2. Institutional setting}\label{sec-institution}

Our setting is a publicly visible, still-contestable opportunity
involving a listed target. A lead buyer is prepared to bid, while
another prospective buyer has not committed to the full diligence and
transaction-preparation process. Target shares trade during that
decision interval. An incumbent is an available acquirer, not the
target's incumbent management; a challenger is a prospective competing
acquirer, not an activist shareholder.

The information relevant to participation can be distributed across
organizations. A buyer may know its own technology, operating
capabilities, and integration capacity, while specialized investors know
the target's customers or product market. Diligence combines these
assessments into a transaction-ready valuation. The private signals in
our extension capture this division: the investor's signal is different
from the buyer's signal, even when the buyer's signal is more accurate.

We distinguish an expression of interest, costly investigation, and a
committed bid. The takeover-process literature establishes the
importance of competition and uncertainty before the visible final
bidding stage (\citeproc{ref-BooneMulherin2007}{Boone and Mulherin
2007}; \citeproc{ref-GentryStroup2019}{Gentry and Stroup 2019}). An
illustrative process record is Imprivata's definitive merger proxy: an
unsolicited approach was followed by deliberation over potential buyers,
confidential outreach, and indications subject to further diligence. The
record also describes concerns about disruption and information leakage
(\citeproc{ref-Imprivata2016}{Imprivata, Inc. 2016}). We use this
example to motivate the separation between interest and costly
participation, not as evidence that stock prices caused entry or that
its confidential phase satisfied our public-visibility condition.

The institutional pilot therefore has a specific task: establish when
the sale opportunity was public, which potential buyer decisions
remained open, and what information prices could add before those
decisions. A stock price existing during confidential negotiations is
not enough. We require an observable interval in which another buyer
could still decide whether to investigate a publicly understood
opportunity.

We isolate target equity as the relevant traded claim. The bidders can
be privately held; alternatively, their securities contain no additional
information about the modeled match. Neither bidder trades the target in
the benchmark. The investor has no control rights relevant to the sale.
The sale procedure binds the whole company, so the mechanism concerns
discovery before participation rather than dispersed-shareholder
tendering. These choices separate our channel from toeholds and takeover
free riding (\citeproc{ref-BulowHuangKlemperer1999}{Bulow et al. 1999};
\citeproc{ref-GrossmanHart1980}{Grossman and Hart 1980}).

\section{3. Related literature}\label{sec-literature}

The closest conceptual comparison is Dow, Goldstein, and Guembel
(\citeproc{ref-DowGoldsteinGuembel2017}{Dow et al. 2017}). Their real
investment decision changes the private incentives to produce
information in financial markets. We share the feedback between real
decisions and information incentives. Our additional object is the
division of acquisition rents between the claim traded by an investor
and the buyer deciding whether to enter. We derive opposite movements in
those returns from an acquisition contest and establish an
entry-deterrence reversal. Information revelation by the investor is
endogenous in our benchmark; the investor's acquisition of its signal is
not.

Edmans, Goldstein, and Jiang
(\citeproc{ref-EdmansGoldsteinJiang2015}{Edmans et al. 2015}) show that
corrective real decisions can discourage adverse-information trading
even under rational pricing. We instead study how rival strength changes
the information sensitivity of target equity and the set of willing
acquirers. Their evidence on the effect of prices on takeovers
(\citeproc{ref-EdmansGoldsteinJiang2012}{Edmans et al. 2012})
establishes the importance of the price-to-control direction, but its
takeover-trigger mechanism is not evidence that a challenger learns its
acquisition match from prices. Luo (\citeproc{ref-Luo2005}{Luo 2005})
studies learning from announcement reactions in completion decisions.
Our participation decision occurs before the buyer set is fixed.

Fishman (\citeproc{ref-Fishman1988}{Fishman 1988}) develops preemptive
bidding that deters costly competition. Hirshleifer and Png
(\citeproc{ref-HirshleiferPng1989}{Hirshleifer and Png 1989}) study
costly investigation, competing bids, and whether facilitating
competition benefits target owners. Gentry and Stroup
(\citeproc{ref-GentryStroup2019}{Gentry and Stroup 2019}) analyze
uncertainty before entry in takeover auctions. We retain the direct
deterrence force and endogenize the separate market information
available before investigation.

Levin and Smith (\citeproc{ref-LevinSmith1994}{Levin and Smith 1994})
characterize auction entry and the coordination costs associated with
potential bidders. Roberts and Sweeting
(\citeproc{ref-RobertsSweeting2013}{Roberts and Sweeting 2013}) show how
costly selective entry affects a seller's choice of sale procedure.
These results make participation an institutional object. We add the
effect of sale payments on information revealed through a distinct
traded claim. Persico (\citeproc{ref-Persico2000}{Persico 2000}) studies
how auction formats shape bidders' incentives to acquire information. In
our mechanism, the relevant financial trader is not an auction
participant, and the claim on which it trades differs from an entrant's
acquisition profit.

Betton, Eckbo, Thompson, and Thorburn
(\citeproc{ref-BettonEtAl2014}{Betton et al. 2014}) analyze merger
negotiations with stock-market feedback and the relation between runups
and offer prices. Lin, Ma, Yang, and Zhu
(\citeproc{ref-LinMaYangZhu2025}{Lin et al. 2025}) model negotiated
payment methods, trading in the merging firms, and subsequent
withdrawal. We study rival discovery and entry under cash consideration
rather than choosing payment composition within an initiated bilateral
deal. Cornelli and Li (\citeproc{ref-CornelliLi2002}{Cornelli and Li
2002}) explain how arbitrageurs' positions affect tendering and generate
an informational advantage; our outside investor instead has information
about a prospective buyer's match and does not determine tendering.

Recent auction-information work sharpens the distinction. Pernoud and
Gleyze (\citeproc{ref-PernoudGleyze2026}{Pernoud and Gleyze 2026}) study
buyers' learning about their own values and competitors. Liu and
Bernhardt (\citeproc{ref-LiuBernhardt2022}{Liu and Bernhardt 2022}) use
post-auction market feedback in security-payment design. Carlin, Liu,
Officer, Pernoud, and Tu (\citeproc{ref-CarlinEtAl2026}{Carlin et al.
2026}) study bidder-pool choice and correlated acquisition values. Our
market operates before participation, and the sale payment determines
which information an outside investor has an incentive to reveal. The
seller-design problem combines these margins rather than taking either
the information experiment or the buyer pool as given.

\section{4. Model}\label{sec-model}

\subsection{4.1 Values, preparation, and the sale
institution}\label{values-preparation-and-the-sale-institution}

We normalize the target's known standalone value to zero and express
acquisition values and preparation costs per target share. Restoring the
same known standalone component to every ownership outcome adds a
constant to prices and payoffs without changing incentives. The informed
order is measured in a small reference trading unit, not as ownership of
the whole target.

The incumbent has private acquisition value \(R\sim U[0,r]\). The
distribution is conditional on public information at the start of the
sale process. Increasing \(r\) is a first-order stochastic
strengthening. The incumbent learns its realization before bidding and
incurs no further participation cost at this stage.

The challenger has acquisition value \(\theta\in\{\ell,h\}\) with equal
prior probabilities. We impose

\[
0<p<\ell<r<h.
\tag{1}
\]

The challenger initially does not know \(\theta\). After observing the
target price it privately learns a preparation cost \(C\), equal to
\(c_L\) with probability \(\rho\) and \(c_H\) otherwise, where
\(0<\rho<1\) and \(0\le c_L<c_H\). Paying the cost reveals its
acquisition value and permits participation. Declining means absence.
Preparation is sunk before bidding. Incumbent value, challenger value,
preparation cost, and noise demand are mutually independent.

The seller commits publicly before trading to a cash second-price
auction with reserve \(p\). The highest admissible bidder acquires the
target and pays the larger of the reserve and the highest competing bid.
Without an admissible bid there is no sale. We implement truthful,
weakly dominant bidding. Conditional on the competing bid, a bidder
wants to win exactly when its value covers the payment required to win.
Its own bid cannot lower that payment conditional on winning. A bid
equal to the reserve is admissible; acquisition-value ties have no
effect in the interior benchmark. Entry at an exactly zero net
preparation payoff is prescribed, a convention that matters at a
posterior plateau.

\subsection{4.2 Trading, information, and
timing}\label{trading-information-and-timing}

The investor observes \(\theta\) and chooses \(q\in[-1,1]\). It has no
initial position, cannot bid for the target, and pays \(k|q|\), where
\(k>0\). Its profit is

\[
q\{V_T-P(X)\}-k|q|,
\qquad X=q+Z,
\qquad f(z)=\frac{1}{2b}e^{-|z|/b},\quad b>1.
\tag{2}
\]

Here \(V_T\) is the terminal target-share payoff and \(Z\) is
independent noise demand. Competitive market makers observe \(X\) and
set

\[
P(X)=\mathbb E[V_T\mid X],
\tag{3}
\]

anticipating preparation and the auction. The challenger observes \(P\),
not \(X\).

The sale opportunity and rule are announced first. The investor then
learns and trades. Market makers price aggregate demand. The challenger
observes the price and cost, chooses preparation, and learns its value
if it prepares. The bidders submit truthful bids, ownership changes if
an admissible bid exists, and financial payoffs are realized.

\subsection{4.3 Equilibrium and
notation}\label{equilibrium-and-notation}

An equilibrium comprises conditional probability distributions
\(\sigma_H,\sigma_L\) over orders, a measurable price function, Bayesian
beliefs based on the observed price, optimal preparation for each cost
realization, and truthful auction bidding. The investor optimizes over
the entire order interval. A unilateral deviation holds the equilibrium
pricing and preparation schedules fixed while changing the distribution
of order flow reaching them. A comparison across values of \(r\) instead
resolves the entire equilibrium.

We use \(t_0\) for expected target proceeds without entry, \(t_H,t_L\)
for proceeds conditional on entry and challenger quality, and
\(g_H,g_L\) for gross challenger acquisition profits. The payoff spread
is \(\Delta_T=t_H-t_L\), and \(B_r(\mu)\) is gross challenger profit at
posterior \(\mu\). The posterior bounds are \(m,M\); the expensive-entry
threshold is \(\tau\) and its order-flow counterpart is \(x^*\).
Favorable-flow probabilities conditional on quality are
\(\alpha_H,\alpha_L\). We denote total entry by \(\mathsf E\),
high-quality ownership by \(\mathsf O_H\), and target revenue by
\(\mathcal R_T\). The boundaries \(r_N,r_U,r_C\) distinguish pooling
existence, sufficient full-order uniqueness, and expensive-entry
feasibility.

\section{5. Acquisition payoffs and the returns to
information}\label{sec-payoffs}

\subsection{5.1 The payment identities}\label{the-payment-identities}

A high-value challenger wins and pays \(\max\{p,R\}\). With a low-value
challenger, the winner pays \(\max\{p,\min(R,\ell)\}\). Without a
challenger, the incumbent pays the reserve exactly when it meets that
reserve. Integration gives

\[
\begin{aligned}
t_0&=p\left(1-\frac p r\right),&
t_H&=\frac r2+\frac{p^2}{2r},&
t_L&=\ell-\frac{\ell^2-p^2}{2r},\\
g_H&=h-\frac r2-\frac{p^2}{2r},&
g_L&=\frac{\ell^2-p^2}{2r},&
\Delta_T(r)&=\frac{(r-\ell)^2}{2r}.
\end{aligned}
\tag{4}
\]

At posterior \(\mu\), the challenger's gross profit and the relevant
derivatives are

\[
\begin{aligned}
B_r(\mu)&=g_L+\mu(g_H-g_L),\\
\Delta_T'(r)&=\frac12-\frac{\ell^2}{2r^2}>0,\\
g_H'(r)&=-\frac12+\frac{p^2}{2r^2}<0,
\qquad g_L'(r)=-\frac{\ell^2-p^2}{2r^2}<0.
\end{aligned}
\tag{5}
\]

Thus \(\partial B_r(\mu)/\partial r<0\) for every fixed posterior. The
decline in acquisition profitability and the increase in target-payoff
sensitivity are different implications of the same sale rule.

\subsection{5.2 General incumbent distributions}\label{prop-payoffs}

\textbf{Proposition 1 (analytical).} \emph{Let the incumbent have a
continuous distribution \(F\) on \([0,\bar r]\), with
\(0<p<\ell<\bar r<h\). A first-order stochastic strengthening of \(F\)
weakly increases the high--low spread of target proceeds and weakly
decreases the challenger's gross acquisition profit at every fixed
posterior. The comparisons are strict when the distribution change has
positive integral over the corresponding ranges below:}

\[
\Delta_T(F)=\mathbb E_F[(R-\ell)_+]
=\int_\ell^{\bar r}[1-F(u)]\,du,
\qquad
G_\theta(F)=\mathbb E_F[(\theta-\max\{p,R\})_+]
=\int_p^\theta F(u)\,du.
\tag{6}
\]

\emph{We extend \(F\) by unity above its support.}

The proof is the payment identity \(T_H(R)-T_L(R)=(R-\ell)_+\). Stronger
competition puts more weight where a high challenger, rather than a low
challenger, changes the price paid to the target. The profit integrand
moves in the opposite direction. Appendix \hyperref[pa-payoffs]{A.1}
supplies the case split and integral argument.

\setcounter{figure}{1}\begin{figure}[H]\centering\includegraphics[width=\linewidth]{figures/two_returns.pdf}
\caption{Incumbent strength raises the sensitivity of target proceeds while reducing challenger acquisition profit at every displayed belief. Panel (a): $\Delta_T(r)$. Panel (b): $B_r(\mu)$ for $\mu\in\{m,1/2,M\}$. All other primitives at the benchmark specification; values per target share.}\label{fig:two_returns}\end{figure}

\section{6. Inference, participation, and residual trading
profits}\label{sec-inference}

\subsection{6.1 Posterior bounds before imposing a trading
profile}\label{lemma-beliefs}

\textbf{Lemma 1 (analytical).} \emph{For arbitrary mixed
state-contingent orders on \([-1,1]\), the order-flow posterior and the
posterior based on price lie in \([m,M]\), where}

\[
\begin{aligned}
a_H(x)&=\int f(x-q)\,d\sigma_H(q),\quad
a_L(x)=\int f(x-q)\,d\sigma_L(q),\\
\mu_X(x)&=\frac{a_H(x)}{a_H(x)+a_L(x)},\qquad
m=\frac1{1+e^{2/b}},\quad M=1-m.
\end{aligned}
\tag{7}
\]

\textbf{Proof.} For any feasible \(q,q'\), the triangle inequality gives

\[
e^{-2/b}\le\frac{f(x-q)}{f(x-q')}\le e^{2/b}.
\tag{8}
\]

Integrating against \(d\sigma_H(q)d\sigma_L(q')\) preserves both
inequalities. Equal priors yield the bounds for \(\mu_X\). Because \(P\)
is a function of \(X\), the price posterior is
\(\mathbb E[\mu_X\mid P]\) and lies in the same interval. \(\square\)

\subsection{6.2 Price sufficiency and optimal
preparation}\label{lemma-price}

At a price-based posterior \(\mu\), preparation is optimal exactly when
\(C\le B_r(\mu)\). Whenever \(c_L<B_r(m)\), entry probability at every
feasible price is at least \(\rho\).

\textbf{Lemma 2 (analytical).} \emph{Suppose \(c_L<B_r(m)\). The
observed price reveals \(\mu_X\) almost surely. Entry and competitive
pricing can be written as}

\[
\begin{aligned}
e_r(\mu)&=\rho+(1-\rho)\mathbf1\{B_r(\mu)\ge c_H\},\\
P_r(\mu)&=t_0+e_r(\mu)[t_L-t_0+\Delta_T\mu].
\end{aligned}
\tag{9}
\]

\emph{The residual advantage of an informed buyer of shares in state
\(H\), and of an informed short seller in state \(L\), is}

\[
\begin{aligned}
A_H(x)&=\mathbb E[V_T\mid H,x]-P(x)=e_r(\mu_X(x))\Delta_T[1-\mu_X(x)],\\
A_L(x)&=P(x)-\mathbb E[V_T\mid L,x]=e_r(\mu_X(x))\Delta_T\mu_X(x),\\
\rho m\Delta_T&\le A_H(x),A_L(x)\le\Delta_T.
\end{aligned}
\tag{10}
\]

\textbf{Proof.} Write \(e(P)\) for entry averaged over the independent
preparation cost. Conditional pricing implies

\[
P=t_0+e(P)[t_L-t_0+\Delta_T\mu_X],
\qquad
\mu_X=\frac{P-t_0-e(P)(t_L-t_0)}{e(P)\Delta_T}.
\tag{11}
\]

The denominator is positive, so the posterior is a measurable function
of the price. Conditioning it again on price leaves it unchanged. The
optimal preparation rule therefore gives (9). For construction,
\(P_r(\mu)\) is strictly increasing: the entry probability is
nondecreasing and positive, while the bracket is positive and strictly
increasing because \(t_L\ge p>t_0\). Subtract the price from
\(t_0+e_r(\mu)(t_H-t_0)\) in state \(H\), and reverse the subtraction in
state \(L\), to obtain (10). \(\square\)

This identity is the core feedback calculation. A publicly anticipated
increase in acquisition proceeds cancels from the informed trader's
profit. The remaining profit is the state-dependent component that noise
trading prevents the market maker from fully identifying.

\section{7. Competition creates competition}\label{sec-results}

\subsection{7.1 The entry-reversal theorem}\label{thm-entry}

\textbf{Theorem 1 (analytical).} \emph{Fix \(0<p<\ell<r_0<r_1<h\),
\(0<\rho<1\), \(b>1\), and \(k>0\). Suppose}

\[
0\le c_L<B_{r_1}(m),
\tag{A1}
\]

\[
B_{r_0}(1/2)<c_H<B_{r_1}(M),
\tag{A2}
\]

\[
\Delta_T(r_0)<k<\left(1-\frac1b\right)\rho m\Delta_T(r_1).
\tag{A3}
\]

\emph{At \(r_0\), the unique equilibrium trading outcome is
\(q_H=q_L=0\), prices are uninformative, and entry is \(\rho\). At
\(r_1\), the unique equilibrium trading outcome is \((q_H,q_L)=(1,-1)\),
the price is informative, and entry strictly exceeds \(\rho\). The
probability that a high-value challenger acquires the target strictly
increases. The strong-incumbent price experiment strictly Blackwell
dominates the weak-incumbent experiment at the same noise law. These
comparisons hold on a nonempty open set of primitives. Uniqueness
concerns trading and on-path entry under truthful auction implementation
and permits arbitrary mixed orders and continuous deviations.}

The conditions describe a region with inexpensive opportunities that are
always worth investigating, expensive opportunities that need favorable
information, and a trading wedge between the weak and strong economies'
information incentives. None assumes that rivalry directly increases
acquisition profit; equation (5) establishes the opposite.

The proof has a short economic sequence. The likelihood-ratio bound
first controls beliefs under every possible trading strategy. The
inexpensive preparation opportunity then ensures a positive entry floor.
That floor makes the observed price sufficient and leaves a positive
residual information advantage. At weak strength, even the largest
possible advantage cannot cover the trading wedge. At strong strength, a
global marginal-profit bound makes increasing the correctly signed order
profitable throughout the order interval. Finally, the informative price
crosses the expensive preparation threshold with positive probability.

For the central bound, fix a candidate equilibrium's residual schedules
and set

\[
F_H(s)=\int f(x-s)A_H(x)\,dx,
\qquad F_L(s)=\int f(x+s)A_L(x)\,dx,
\qquad U_\theta(s)=sF_\theta(s)-ks.
\tag{12}
\]

The density satisfies \(|f'|\le f/b\), giving
\(|F_\theta'|\le F_\theta/b\). Consequently,

\[
U_\theta'(s)\ge\left(1-\frac{s}{b}\right)F_\theta(s)-k
\ge\left(1-\frac1b\right)\rho m\Delta_T-k>0
\tag{13}
\]

almost everywhere in the strong economy. This controls every continuous
deviation, not merely candidate first-order conditions. Appendix
\hyperref[pa-entry]{A.2} supplies each logical step, including
equilibrium construction and the information ordering.

Under full orders the expensive buyer enters for \(X\ge x^*\), where

\[
\begin{aligned}
\tau&=\frac{c_H-g_L}{g_H-g_L},&
x^*&=\frac b2\log\frac\tau{1-\tau},\\
\alpha_H&=1-\frac12e^{(x^*-1)/b},&
\alpha_L&=\frac12e^{-(x^*+1)/b},\\
\mathsf E&=\rho+\frac{1-\rho}{2}(\alpha_H+\alpha_L),&
\mathsf O_H&=\frac12[\rho+(1-\rho)\alpha_H].
\end{aligned}
\tag{14}
\]

Condition (A2) places \(\tau\in(1/2,M)\) and \(x^*\in(0,1)\).

\subsection{7.2 Fixed information and exact
boundaries}\label{prop-fixed}

\textbf{Proposition R1 (analytical).} \emph{Hold the joint distribution
of a signal, challenger quality, and preparation cost fixed as \(r\)
changes. If the posterior generated by that signal does not change with
\(r\), entry is weakly decreasing in \(r\). It decreases strictly when
the decline in gross profit crosses preparation costs on a
positive-probability set of signal--cost realizations.}

The proof is pointwise: \(\mathbf1\{C\le B_r(\mu)\}\) cannot switch from
nonentry to entry as \(B_r(\mu)\) falls. This applies within a region of
fixed full orders. It does not apply merely because prices are
informative: changing orders changes the information experiment.

\subsection{7.3 Pooling, full orders, and expensive-entry
feasibility}\label{prop-thresholds}

\textbf{Proposition 3 (analytical).} \emph{Restrict attention to
strengths in \((\ell,h)\) for which \(c_L<B_r(m)\) and \(B_r(1/2)<c_H\).
Define}

\[
\begin{aligned}
\mathfrak r(d)&=\ell+d+\sqrt{d^2+2\ell d},\\
r_N&=\mathfrak r(2k/\rho),\qquad
r_U=\mathfrak r\left(\frac{k}{(1-1/b)\rho m}\right).
\end{aligned}
\tag{15}
\]

\emph{Pooling exists exactly when \(\rho\Delta_T(r)/2\le k\),
equivalently \(r\le r_N\). The stronger restriction \(r<\mathfrak r(k)\)
guarantees its uniqueness. The restriction \(r>r_U\) guarantees unique
full orders. Expensive entry is impossible in every equilibrium when
\(B_r(M)<c_H\). If the equality \(B_r(M)=c_H\) has a solution in the
specified domain, its unique crossing is \(r_C\), given by the relevant
root}

\[
r_C=\frac{Mh-c_H+\sqrt{(Mh-c_H)^2+M[(1-M)\ell^2-p^2]}}{M}.
\tag{16}
\]

\emph{At \(r_C\), the prescribed tie rule is retained.}

The exact pooling-existence boundary, a sufficient uniqueness boundary,
and the expensive-entry feasibility boundary answer different questions.
In particular, pooling existence does not exclude an informative
equilibrium at the same strength. Under Laplace noise the maximum
posterior occurs on a positive-probability tail, so indifference at
\(r_C\) cannot be discarded as a null event. Appendix
\hyperref[pa-thresholds]{A.3} derives the formulas and specifies their
domain.

\subsection{7.4 A rise and fall across uniquely determined
economies}\label{cor-three}

\textbf{Corollary 3 (analytical).} \emph{Let \(r_0,r_1\) satisfy Theorem
1 and choose \(r_2\in(r_1,h)\) such that \(c_L<B_{r_2}(m)\),
\(B_{r_2}(M)<c_H\), and \(k<(1-1/b)\rho m\Delta_T(r_2)\). Then each
economy has a unique trading and entry outcome and}

\[
\mathsf E(r_0)=\rho<\mathsf E(r_1),
\qquad \mathsf E(r_2)=\rho.
\tag{17}
\]

We implement this comparison at \(r_0=1.2\), \(r_1=3\), and \(r_2=3.6\).
Entry is respectively 0.250000, 0.522757, and 0.250000. This finite
nonmonotonicity result does not assign a unique outcome throughout the
intermediate correspondence.

\subsection{7.5 Asymmetric informative equilibria}\label{prop-certified}

\textbf{Proposition 4 (computer-assisted).} \emph{At the benchmark
primitives specified in Appendix \hyperref[pa-parameters]{A.9}, there
exist informative equilibria \((q_H,q_L)=(1,-v_j)\) at strengths \(r_j\)
whose certified enclosures are}

\[
\begin{aligned}
r_{\mathrm a}&=1.55,&v_{\mathrm a}&\in[0.46031618,\,0.46031620],&\mathsf E_{\mathrm a}&\in[0.5450528898,\,0.5450528922],\\
r_{\mathrm b}&=1.60,&v_{\mathrm b}&\in[0.70747537,\,0.70747539],&\mathsf E_{\mathrm b}&\in[0.5487563062,\,0.5487563085],\\
r_{\mathrm c}&=1.65,&v_{\mathrm c}&\in[0.90333198,\,0.90333201],&\mathsf E_{\mathrm c}&\in[0.5513607988,\,0.5513608020].
\end{aligned}
\tag{18}
\]

\emph{The entry intervals are strictly ordered upward. Each economy also
admits pooling with entry \(\rho\).}

The investor buys maximally after favorable information but chooses an
interior short after unfavorable information. The latter problem is
globally strictly concave. Interval sign changes establish exact
interior best responses, and interval bounds on the favorable type's
marginal profit exclude all smaller purchases. We therefore certify
continuous-action equilibria rather than equating a small residual with
equilibrium. Appendix \hyperref[pa-certificate]{A.4} contains the proof
steps and the inequalities completing the certificates. The online
appendix gives the exact integration formulas and arithmetic protocol.

\setcounter{figure}{0}\begin{figure}[H]\centering\includegraphics[width=\linewidth]{figures/equilibrium_correspondence.pdf}
\caption{The equilibrium correspondence at the benchmark primitives. Panel (a): total entry $\mathsf E$ against incumbent strength $r$ for every accepted branch: pooling (no trade), full orders, the asymmetric family $(q_H,q_L)=(1,-v)$ by numerical continuation, and any other pure or mixed profiles found. Shaded regions are analytically established uniqueness regions (no trade below $\mathfrak r(k)$; full orders above $r_U$). Dashed verticals mark $r_N$ (exact pooling-existence boundary), $r_U$ (sufficient full-order uniqueness boundary), and $r_C$ (expensive-entry feasibility boundary). Certified points carry interval enclosures. Lines are broken at unresolved nodes or branch changes; a missing branch is not a uniqueness label. Panel (b): the unfavorable-state order magnitude $v$ along asymmetric candidates, with full orders as the boundary.}\label{fig:correspondence}\end{figure}

\section{8. Information and institutions beyond the
benchmark}\label{sec-extensions}

\subsection{8.1 Smooth noise and atomless
preparation}\label{cor-logistic}

\textbf{Corollary 1 (analytical).} \emph{Replace Laplace noise by
logistic noise with density}

\[
f_{\mathrm{log}}(z)=\frac{1}{4b}\operatorname{sech}^{2}\left(\frac z{2b}\right),\qquad b>1.
\tag{19}
\]

\emph{Under (A1)--(A3), the unique trading outcomes and the entry and
ownership comparisons of Theorem 1 remain valid.}

The log-density derivative is bounded in absolute value by \(1/b\),
which preserves the global trading argument. Under full orders, the
posterior is strictly increasing and ranges over \((m,M)\). It
approaches the bounds smoothly rather than becoming constant in the
tails. If \(A=e^{1/b}\) and \(w=\sqrt{\tau/(1-\tau)}\), the entry
threshold and favorable-flow probabilities are

\[
x^*_{\mathrm{log}}=b\log\frac{Aw-1}{A-w},\qquad
\alpha_H^{\mathrm{log}}=\frac1{1+e^{(x^*_{\mathrm{log}}-1)/b}},\qquad
\alpha_L^{\mathrm{log}}=\frac1{1+e^{(x^*_{\mathrm{log}}+1)/b}}.
\tag{20}
\]

\subsection{8.2 Preparation-cost heterogeneity}\label{cor-costs}

\textbf{Corollary 2 (analytical).} \emph{Replace the cost atoms by
atomless low- and high-cost distributions with probabilities \(\rho\)
and \(1-\rho\), supported respectively within
\([c_L-\varepsilon_C,c_L+\varepsilon_C]\) and
\([c_H-\varepsilon_C,c_H+\varepsilon_C]\). Suppose}

\[
\begin{gathered}
c_L-\varepsilon_C\ge0,\qquad c_L+\varepsilon_C<B_{r_1}(m),\\
B_{r_0}(1/2)<c_H-\varepsilon_C<c_H+\varepsilon_C<B_{r_1}(M),
\end{gathered}
\tag{21}
\]

\emph{and retain (A3). Theorem 1's conclusions hold under either noise
specification.}

Every low-cost realization participates at every feasible belief.
High-cost realizations stay out at the weak prior and participate after
sufficiently favorable strong-economy prices. The price is still
strictly increasing in the posterior because the preparation-cost
distribution generates a nondecreasing entry rule with a positive floor.
Appendix \hyperref[pa-smooth]{A.5} gives both corollary proofs.

The effect is not confined to widely separated acquisition values. The
moderate-value specification has \(h=2\), \(\ell=1\), and satisfies each
strict inequality of Theorem 1. Entry rises from 0.250000 to 0.526805.
Appendix \hyperref[pa-smooth]{A.5} also constructs a nonempty region for
every \(h>\ell\) by choosing the strengths near \(\ell\) and a
compatible trading wedge.

\subsection{8.3 Complementary private information}\label{thm-signals}

We now give both the trader and the buyer imperfect private information.
Let \(T,Y\in\{+,-\}\) be conditionally independent given challenger
quality, with

\[
\Pr(T=+\mid H)=\Pr(T=-\mid L)=a,\qquad
\Pr(Y=+\mid H)=\Pr(Y=-\mid L)=d,
\qquad a,d\in(1/2,1).
\tag{22}
\]

The trader observes \(T\) but not \(Y\). The buyer observes \(Y\) and
the target price before preparation. The buyer's signal may be more
accurate: \(d>a\) is allowed. Diligence still reveals the acquisition
value before bidding. All remaining independence and timing assumptions
are unchanged.

Let \(\lambda_X=\Pr(T=+\mid X)\). Public and joint buyer posteriors are

\[
\begin{aligned}
\mu_X&=(1-a)+(2a-1)\lambda_X,\quad
\mu_-=(1-a)+(2a-1)m,\quad \mu_+=(1-a)+(2a-1)M,\\
\phi_+(\mu)&=\frac{d\mu}{d\mu+(1-d)(1-\mu)},\qquad
\phi_-(\mu)=\frac{(1-d)\mu}{(1-d)\mu+d(1-\mu)}.
\end{aligned}
\tag{23}
\]

Write \(w_H=t_H-t_0\) and \(w_L=t_L-t_0\). The buyer's private signal
changes entry conditional on fundamental quality. With
\(I_y=\mathbf1\{B_r(\phi_y(\mu))\ge c_H\}\),

\[
\begin{aligned}
e_H&=\rho+(1-\rho)[dI_++(1-d)I_-],\\
e_L&=\rho+(1-\rho)[(1-d)I_++dI_-],\\
D&=e_Hw_H-e_Lw_L,\qquad
\rho\Delta_T\le D\le\Delta_T+(1-\rho)(2d-1)w_H.
\end{aligned}
\tag{24}
\]

We cannot replace these conditional probabilities by a common entry
rate. Competitive pricing instead gives

\[
P=t_0+e_L(P)w_L+\mu_XD(P),\qquad
\mu_X=\frac{P-t_0-e_L(P)w_L}{D(P)}.
\tag{25}
\]

The price reveals the public posterior; the buyer combines it with its
private information. The favorable and unfavorable trader-signal
residuals become

\[
A_+=(2a-1)(1-\lambda_X)D,
\qquad A_-=(2a-1)\lambda_XD.
\tag{26}
\]

\textbf{Theorem R2 (analytical).} \emph{In the complementary-signal
economy, choose \(0<p<\ell<r_0<r_1<h\) and suppose}

\[
\begin{gathered}
c_L<B_{r_1}(\phi_-(\mu_-)),\qquad
B_{r_0}(d)<c_H<B_{r_1}(\phi_+(\mu_+)),\\
(2a-1)\{\Delta_T(r_0)+(1-\rho)(2d-1)w_H(r_0)\}<k\\
<\left(1-\frac1b\right)m(2a-1)\rho\Delta_T(r_1).
\end{gathered}
\tag{27}
\]

\emph{The weak economy has unique zero informed orders and entry
\(\rho\). The strong economy has unique orders \(q(T=+)=1\) and
\(q(T=-)=-1\), and entry strictly exceeds \(\rho\). The conclusions
permit arbitrary mixed signal-contingent orders and every continuous
unilateral deviation.}

The proof controls public beliefs before imposing any order profile,
proves the price inversion with state-dependent entry, and applies
global trading bounds to (26). In the weak economy, even the buyer's
favorable private signal is insufficient for expensive preparation. In
the strong economy, that private signal combined with a favorable price
makes preparation worthwhile. Appendix \hyperref[pa-signals]{A.6} states
every step.

Our example uses buyer accuracy 75\% and trader accuracy 70\%. Entry
rises from 0.850000 to 0.879438. The mechanism uses complementary
information, not an investor whose information uniformly dominates that
of the buyer.

\subsection{8.4 What payment property supports
discovery?}\label{lemma-bargaining}

To isolate the role of payment design, we consider a separate
verifiable-value bargaining institution. We set the reserve to zero.
After diligence, values are publicly verifiable and the highest-value
buyer acquires the target. A sale to the next-best buyer at its value is
an enforceable fallback, with zero-surplus acceptance prescribed. The
seller receives share \(\eta\) of the winning buyer's surplus above that
fallback in Nash bargaining. Without a challenger, its fallback is zero
and it receives \(\eta R\).

\textbf{Lemma 3 (analytical).} \emph{Let \(R\) have any distribution
supported on \([0,R_{\max}]\), where \(0<\ell<R_{\max}<h\). For
\(0\le\eta<1\), this institution generates}

\[
\begin{aligned}
T_\eta(R,\theta)&=(1-\eta)\min\{R,\theta\}+\eta\max\{R,\theta\},\\
G_{\theta,\eta}&=(1-\eta)\mathbb E[(\theta-R)_+],\\
\Delta_\eta&=\eta(h-\ell)+(1-2\eta)\mathbb E[(R-\ell)_+].
\end{aligned}
\tag{28}
\]

\emph{A first-order stochastic strengthening of the incumbent weakly
reduces challenger profits. It weakly increases the target's information
spread for \(\eta<1/2\), leaves that spread unchanged at \(\eta=1/2\),
and weakly decreases it for \(\eta>1/2\). Strictness follows from a
positive integral change in the corresponding payoff function.}

Runner-up discipline, rather than the auction label alone, drives the
original opposition. When the seller already captures much of the
winner's own value, stronger competition need not make the target claim
more sensitive to challenger quality. This result characterizes the
acquisition stage of a specified alternative institution; solving its
complete trading and entry game is part of the sale-design program. The
case split is in Appendix \hyperref[pa-bargaining]{A.7}.

\setcounter{figure}{3}\begin{figure}[H]\centering\includegraphics[width=\linewidth]{figures/bargaining_weight.pdf}
\caption{Bargaining and the division of information-sensitive returns in the verifiable-value institution with zero reserve, benchmark $h=10$, $\ell=1$, and uniform incumbents with $r=1.2$ (weak) and $r=3$ (strong). Panel (a): $\Delta_\eta$ against the seller weight $\eta$; the strength ordering reverses at $\eta=1/2$. Panel (b): challenger profits $G_{H,\eta}$ and $G_{L,\eta}$ (log scale). These are acquisition-stage comparisons, not equilibrium entry predictions.}\label{fig:bargaining}\end{figure}

\section{9. The real value of discovery}\label{sec-welfare}

\subsection{9.1 Access to prices and acquisition
surplus}\label{prop-welfare}

\textbf{Proposition 2 (analytical).} \emph{Hold \(r=r_1\) and all other
primitives fixed under Theorem 1. Compare the feedback equilibrium with
the equilibrium in which the challenger cannot observe the target price,
reoptimizing trading and pricing in both. Access to prices strictly
increases expected target proceeds and acquisition surplus net of
preparation costs. The comparison also holds under Corollaries 1 and 2.}

The no-feedback buyer enters only at low cost. The same global trading
bound supports full orders in both economies, so the financial
information experiment and real trading costs are matched. What changes
is the buyer's access to that experiment.

Without entry the incremental ownership value is
\(W_0(R)=R\mathbf1\{R\ge p\}\). With a challenger of value \(\theta>p\)
it is \(W_\theta(R)=\max\{R,\theta\}\). The key identity is

\[
W_\theta(R)-W_0(R)
=(\theta-\max\{p,R\})_++p\mathbf1\{R<p\}.
\tag{29}
\]

Thus an additional preparation decision at posterior \(\mu\) and cost
\(C\) produces conditional expected net surplus \(B_r(\mu)-C+p^2/r\).
Whenever the buyer chooses that additional entry, its first term is
nonnegative and the second is positive. Low-cost entry is unchanged. For
cost atoms,

\[
\begin{aligned}
\Delta\mathcal W&=(1-\rho)\mathbb E\left[
\left(B_r(\mu_X)-c_H+\frac{p^2}{r}\right)\mathbf1\{\mu_X\ge\tau\}\right]>0,\\
\Delta\mathcal R_T&=\frac{1-\rho}{2}
\left[\alpha_H(t_H-t_0)+\alpha_L(t_L-t_0)\right]>0.
\end{aligned}
\tag{30}
\]

Payments between bidders and target owners are transfers in this surplus
calculation. Appendix \hyperref[pa-welfare]{A.8} completes the
comparison and integrates over atomless costs. The result concerns
access to prices at fixed incumbent strength, not a welfare ordering
over all changes in competition or sale rules.

\subsection{9.2 Separating information from price
levels}\label{separating-information-from-price-levels}

At strong incumbent strength, mean target proceeds are 0.872392 with
price feedback and 0.614583 when the buyer cannot observe the price. We
attach a deterministic external dividend of 0.257809 to the traded claim
in the latter economy. This matches the mean financial payoff and price
while leaving every trading residual and entry decision unchanged. The
dividend is a level-matching diagnostic, not a sale term available to
the seller or a resource gain in the welfare calculation.

\section{10. Numerical illustration}\label{sec-numerics}

We quantify the mechanism using a benchmark, a moderate-value economy,
and separate extensions. Each exercise begins with its own primitive
vector and resolves the relevant equilibrium; the different extensions
are not treated as a single jointly solved economy. The full input
ledger, acceptance checks, and quantity definitions are in Online
Appendix C.

In the benchmark, strengthening the incumbent lowers prior gross
challenger profit from 4.804167 to 4.291667, while increasing the
target's information spread from 0.016667 to 0.666667. Entry
nevertheless rises from 0.250000 to 0.522757. The probability of
high-quality ownership rises from 0.125000 to 0.324192.

\setcounter{table}{0}\input{tables/table1_auction_primitives.tex}

The frozen-information control imposes the same informative order
profile at both strengths while allowing the buyer to reoptimize. Its
entry decreases from 0.562178 to 0.522757. Removing access to prices
yields entry 0.250000 and 0.250000. These controls distinguish the
direct deterrence effect from the endogenous supply and use of price
information. The frozen informative profile is not an equilibrium in the
weak economy.

\setcounter{table}{1}\input{tables/table2_equilibrium_controls.tex}

Under logistic noise, strong-economy entry is 0.301509, compared with
0.522757 under Laplace noise. The order-flow threshold is 5.424598398,
and its distance in noise standard deviations is 1.495369. Both
experiments have the same posterior bounds at the same scale parameter,
but they place different mass near those bounds. We therefore compare
upper-tail probabilities directly rather than treating scale or variance
as an information ordering.

\setcounter{figure}{2}\begin{figure}[H]\centering\includegraphics[width=\linewidth]{figures/posterior_tail_entry.pdf}
\caption{The upper tail of price information under full orders at the strong benchmark strength and scale $b=2$. Panel (a): $\Pr(\mu_X\ge\tau)$ against the threshold distance $M-\tau$ for Laplace and logistic noise. Panel (b): implied total entry $\rho+(1-\rho)\Pr(\mu_X\ge\tau)$ and the conditional favorable-flow probabilities $\alpha_H,\alpha_L$. At $M-\tau=0$ the Laplace plateau enters under the tie rule (dot) while the logistic mass is zero (cross). These are fixed-profile information diagnostics; the accompanying data file records for each threshold whether the equilibrium interpretation with the implied high cost is validated.}\label{fig:posterior_tail}\end{figure}

Atomless preparation costs yield strong-economy entry 0.522715 with
Laplace noise and 0.301374 with logistic noise. The moderate-value and
complementary-signal cases retain positive margins in their respective
analytical conditions.

\setcounter{table}{2}\input{tables/table3_extensions.tex}

\section{11. Sale design as information policy}\label{sec-design}

\subsection{11.1 The seller's continuation problem}\label{open-design}

\textbf{Research question A (open).} \emph{Which sale rules maximize
target proceeds when they change the information revealed before bidder
participation? We seek a characterization of seller-optimal discovery,
its dependence on incumbent strength, and the conditions under which it
reverses or reinforces entry deterrence.}

We begin with a reserve chosen publicly before trading. For every
reserve \(p\), let \(\mathcal E(p,r)\) be the set of trading, pricing,
and entry continuations. If \(\sigma\in\mathcal E(p,r)\), write
\(e_H(p,r;\sigma),e_L(p,r;\sigma)\) for state-specific entry. Seller
revenue is

\[
\mathcal R_T(p,r;\sigma)
=t_0(p,r)+\frac12\sum_{\theta\in\{H,L\}}e_\theta(p,r;\sigma)
[t_\theta(p,r)-t_0(p,r)].
\tag{31}
\]

A seller equilibrium specifies a continuation after every feasible
reserve, not only after its chosen reserve. Given such a selection
\(\sigma^*(p)\in\mathcal E(p,r)\), the reserve must satisfy

\[
p^*\in\arg\max_{p\in[0,h]}\mathcal R_T(p,r;\sigma^*(p)).
\tag{32}
\]

Existence, attainment, and continuation selection are part of the result
to establish. Optimistic and pessimistic revenue envelopes describe the
retained correspondence; neither envelope is automatically an
equilibrium objective. The strongest comparative-static target is an
entry reversal under seller-optimal terms. A characterization of when an
optimal seller instead eliminates the reversal would answer the same
design question.

\subsection{11.2 Payoffs over the complete reserve
domain}\label{payoffs-over-the-complete-reserve-domain}

The reliable starting point is the realized sale rule. For any realized
challenger value \(v\),

\[
\begin{aligned}
t_0(p,r)&=\mathbb E[p\mathbf1\{R\ge p\}],\\
t_v(p,r)&=
\begin{cases}
\mathbb E[\max\{p,\min(R,v)\}],&v\ge p,\\
t_0(p,r),&v<p,
\end{cases}\\
g_v(p,r)&=\mathbb E[(v-\max\{p,R\})_+].
\end{aligned}
\tag{33}
\]

For \(p<\ell\), equation (4) applies. For \(\ell<p<r\), \(t_L=t_0\),
\(g_L=0\), \(t_H=r/2+p^2/(2r)\), and \(g_H=h-t_H\). For \(r\le p<h\),
\(t_0=t_L=g_L=0\), \(t_H=p\), and \(g_H=h-p\). Equality at an
acquisition-value atom is evaluated using the stated admissibility
convention. A reserve above the highest possible value prevents a sale.

These regimes change the claim traded by investors. Excluding a
low-value buyer can increase the target's high--low payoff spread,
making information valuable even against a weak incumbent. With atomless
acquisition values, we integrate (33) over the conditional value
distribution rather than extending a formula through an exclusion
boundary.

\subsection{11.3 The local extraction--discovery
decomposition}\label{the-local-extractiondiscovery-decomposition}

On a differentiable continuation branch with atomless preparation costs
and \(0<p<\ell\), let \(\mathsf E=(e_H+e_L)/2\). Differentiating seller
revenue gives

\[
\frac{d\mathcal R_T}{dp}
=(1-\mathsf E)\left(1-\frac{2p}{r}\right)+\mathsf E\frac p r
+\frac12\sum_\theta\frac{de_\theta}{dp}(t_\theta-t_0).
\tag{34}
\]

If the cost CDF is \(H_C\) with density \(h_C\), and \(\mu_\theta(z;p)\)
is the posterior reached by fundamental state \(\theta\) at noise
realization \(z\), then

\[
\frac{de_\theta}{dp}
=\int f(z)h_C(B_{p,r}(\mu_\theta))
\left[-\frac p r+(g_H-g_L)\frac{d\mu_\theta(z;p)}{dp}\right]dz.
\tag{35}
\]

The first term inside the bracket is direct rent extraction. The second
is the change in information generated by equilibrium trading. It
vanishes locally on a fixed full-order branch, but not generally when
orders adjust. Equations (34)--(35) are local decompositions, not global
sign restrictions. Appendix \hyperref[pa-design]{A.10} gives the
differentiability conditions and derivation.

\subsection{11.4 A reserve diagnostic with atomless acquisition
values}\label{diagnostic-reserve}

\textbf{Diagnostic 1 (numerical diagnostic).} \emph{Let challenger
quality indicate a low or high value class, with equally likely classes.
Conditional on class, acquisition value is uniform on
\([\ell-\varepsilon_V,\ell+\varepsilon_V]\) or
\([h-\varepsilon_V,h+\varepsilon_V]\). The investor observes the class;
preparation reveals the exact value. At the specifications in Online
Appendix C.6, the tested reserve increase raises seller proceeds at both
strengths and produces informative trading in both economies. The stated
trading outcomes satisfy the corresponding analytical global bounds.}

The value half-width is 0.05. Raising the reserve from 0.5 to 1.1 raises
weak-economy revenue from 0.392665 to 0.432173, and strong-economy
revenue from 0.872367 to 1.014500. Entry under the higher reserve is
0.540309 and 0.511638, respectively. This profitable alternative
identifies a substantive seller-design margin with atomless acquisition
values. Its optimization is the next result, not a property inferred
from this comparison.

\setcounter{table}{3}\input{tables/table4_reserve_comparisons.tex}

\section{12. Research program and empirical design}\label{sec-program}

We pursue the seller-design problem first, retaining the complete
continuation problem at each reserve and making challenger values
continuous. The objective is an institutional result about how a seller
obtains informative participation. The comparison between a fixed
institution and an optimized one determines whether the seller uses,
strengthens, or substitutes for the incumbent-induced information
effect.

Next we characterize the intermediate equilibrium correspondence. The
certified asymmetric points provide exact anchors, not a complete
branch. Continuation must allow changing buy--sell asymmetry and mixed
orders. We seek conditions for branch existence, local continuation, and
multiplicity, followed by a selection argument where one is economically
justified.

The next institutional result concerns commitment. \textbf{Research
question B (open).} \emph{When does fixing sale terms before trading
improve the seller's ex ante proceeds relative to revising them after
observing the price but before preparation?} Anticipated repricing can
change the return to revealing information and the return to becoming an
informed buyer. We will solve that timing rather than impose a
price-indexed penalty with the desired sign.

We then solve an alternative acquisition institution in full. The
bargaining lemma identifies the payment property to study. In a
first-price alternative, bids and expected payments depend on the
incumbent's posterior about a challenger selected through
price-dependent entry; those posteriors belong in the auction
continuation. Endogenous investor research likewise requires solving the
uninformed trader's incentives, including potential manipulation through
real decisions (\citeproc{ref-GoldsteinGuembel2008}{Goldstein and
Guembel 2008}). These are separate games, not substitutions into
unchanged payoff formulas.

The empirical component begins with an institutional pilot. We will
select publicly visible, still-open sale opportunities involving listed
targets and reconstruct the timing of initial approaches, public
visibility, contacts, diligence, substantive proposals, revisions, and
final selection. At each event we record what was publicly known at that
time, separating later retrospective disclosure from contemporaneous
availability. Online Appendix D defines the selection rule, coding
fields, source hierarchy, and adjudication protocol. The pilot is a
design scaffold: it does not report a selected sample, measured effect,
or identified causal instrument.

The measurement priorities follow directly from the theory. Entry is a
decision to investigate or submit a substantive bid, not simply a count
of public offers. Incumbent strength must be measured using information
preceding the challenger's decision; a final winning bid is not an
exogenous initial threat. Financial information must also precede entry.
A relation between returns and subsequent bidders can reflect
anticipation of entry as well as learning from prices. The first
empirical task is therefore to establish the modeled decision interval
and its information sets. A targeted causal or structural exercise
follows the final institutional theorem.

\section{13. Conclusion}\label{sec-conclusion}

Competition changes both acquisition profits and the financial return to
revealing acquisition information. We show that the second effect can
make a stronger incumbent attract another buyer, even though it reduces
that buyer's profit at every fixed belief. The result survives
complementary private information and smooth distributions. The
intermediate correspondence shows why informative markets need not
behave like a fixed information experiment, while the bargaining
comparison identifies the relevant division of acquisition value.

A company is not sold to a fixed list of fully informed buyers. Its sale
rules and stock market help determine who becomes an informed, willing
buyer. \textbf{How should corporate-control institutions be designed
once their effect on discovery is taken seriously?} That is the
seller-design question our equilibrium results make concrete.

\section{References}\label{references}

\protect\phantomsection\label{refs}
\begin{CSLReferences}{1}{1}
\bibitem[\citeproctext]{ref-BettonEtAl2014}
Betton, Sandra, B. Espen Eckbo, Rex Thompson, and Karin S. Thorburn.
2014. {``Merger Negotiations with Stock Market Feedback.''} \emph{The
Journal of Finance} 69 (4): 1705--45.
\url{https://doi.org/10.1111/jofi.12151}.

\bibitem[\citeproctext]{ref-BooneMulherin2007}
Boone, Audra L., and J. Harold Mulherin. 2007. {``How Are Firms Sold?''}
\emph{The Journal of Finance} 62 (2): 847--75.
\url{https://doi.org/10.1111/j.1540-6261.2007.01225.x}.

\bibitem[\citeproctext]{ref-BulowHuangKlemperer1999}
Bulow, Jeremy, Ming Huang, and Paul Klemperer. 1999. {``Toeholds and
Takeovers.''} \emph{Journal of Political Economy} 107 (3): 427--54.
\url{https://doi.org/10.1086/250068}.

\bibitem[\citeproctext]{ref-CarlinEtAl2026}
Carlin, Bruce I., Tingting Liu, Micah S. Officer, Agathe Pernoud, and
Danni Tu. 2026. \emph{Bidder Pools in Mergers and Acquisitions}. Working
Paper No. 34846. National Bureau of Economic Research.
\url{https://doi.org/10.3386/w34846}.

\bibitem[\citeproctext]{ref-CornelliLi2002}
Cornelli, Francesca, and David D. Li. 2002. {``Risk Arbitrage in
Takeovers.''} \emph{The Review of Financial Studies} 15 (3): 837--68.
\url{https://doi.org/10.1093/rfs/15.3.837}.

\bibitem[\citeproctext]{ref-DowGoldsteinGuembel2017}
Dow, James, Itay Goldstein, and Alexander Guembel. 2017. {``Incentives
for Information Production in Markets Where Prices Affect Real
Investment.''} \emph{Journal of the European Economic Association} 15
(4): 877--909. \url{https://doi.org/10.1093/jeea/jvw023}.

\bibitem[\citeproctext]{ref-EdmansGoldsteinJiang2012}
Edmans, Alex, Itay Goldstein, and Wei Jiang. 2012. {``The Real Effects
of Financial Markets: The Impact of Prices on Takeovers.''} \emph{The
Journal of Finance} 67 (3): 933--71.
\url{https://doi.org/10.1111/j.1540-6261.2012.01738.x}.

\bibitem[\citeproctext]{ref-EdmansGoldsteinJiang2015}
Edmans, Alex, Itay Goldstein, and Wei Jiang. 2015. {``Feedback Effects,
Asymmetric Trading, and the Limits to Arbitrage.''} \emph{American
Economic Review} 105 (12): 3766--97.
\url{https://doi.org/10.1257/aer.20141271}.

\bibitem[\citeproctext]{ref-Fishman1988}
Fishman, Michael J. 1988. {``A Theory of Preemptive Takeover Bidding.''}
\emph{The RAND Journal of Economics} 19 (1): 88--101.
\url{https://doi.org/10.2307/2555399}.

\bibitem[\citeproctext]{ref-GentryStroup2019}
Gentry, Matthew, and Caleb Stroup. 2019. {``Entry and Competition in
Takeover Auctions.''} \emph{Journal of Financial Economics} 132 (2):
298--324. \url{https://doi.org/10.1016/j.jfineco.2018.10.007}.

\bibitem[\citeproctext]{ref-GoldsteinGuembel2008}
Goldstein, Itay, and Alexander Guembel. 2008. {``Manipulation and the
Allocational Role of Prices.''} \emph{The Review of Economic Studies} 75
(1): 133--64. \url{https://doi.org/10.1111/j.1467-937X.2007.00467.x}.

\bibitem[\citeproctext]{ref-GrossmanHart1980}
Grossman, Sanford J., and Oliver D. Hart. 1980. {``Takeover Bids, the
Free-Rider Problem, and the Theory of the Corporation.''} \emph{The Bell
Journal of Economics} 11 (1): 42--64.
\url{https://doi.org/10.2307/3003400}.

\bibitem[\citeproctext]{ref-HirshleiferPng1989}
Hirshleifer, David, and I. P. L. Png. 1989. {``Facilitation of Competing
Bids and the Price of a Takeover Target.''} \emph{The Review of
Financial Studies} 2 (4): 587--606.
\url{https://doi.org/10.1093/rfs/2.4.587}.

\bibitem[\citeproctext]{ref-Imprivata2016}
Imprivata, Inc. 2016. \emph{Definitive Merger Proxy Statement, Form
{DEFM14A}}. U.S. Securities and Exchange Commission.
\url{https://www.sec.gov/Archives/edgar/data/1328015/000119312516677939/d226798ddefm14a.htm}.

\bibitem[\citeproctext]{ref-LevinSmith1994}
Levin, Dan, and James L. Smith. 1994. {``Equilibrium in Auctions with
Entry.''} \emph{The American Economic Review} 84 (3): 585--99.
\url{https://www.jstor.org/stable/2118069}.

\bibitem[\citeproctext]{ref-LinMaYangZhu2025}
Lin, Tse-Chun, Xiaorong Ma, Liyan Yang, and Minxing Zhu. 2025.
\emph{Payment Methods and Market Feedback in Mergers and Acquisitions}.
Working Paper 2025/189. HKU Jockey Club Enterprise Sustainability Global
Research Institute. \url{https://doi.org/10.2139/ssrn.5789982}.

\bibitem[\citeproctext]{ref-LiuBernhardt2022}
Liu, Tingjun, and Dan Bernhardt. 2022. {``Simplifying Auction Designs
via Market Feedback.''}
\url{https://warwick.ac.uk/fac/soc/economics/staff/mdbernhardt/market_feedback_2022res_final.pdf}.

\bibitem[\citeproctext]{ref-Luo2005}
Luo, Yuanzhi. 2005. {``Do Insiders Learn from Outsiders? Evidence from
Mergers and Acquisitions.''} \emph{The Journal of Finance} 60 (4):
1951--82. \url{https://doi.org/10.1111/j.1540-6261.2005.00784.x}.

\bibitem[\citeproctext]{ref-PernoudGleyze2026}
Pernoud, Agathe, and Simon Gleyze. 2026. {``How Competition Shapes
Information in Auctions.''}
\url{https://agathepernoud.com/Pernoud_Gleyze_InfoCompetition.pdf}.

\bibitem[\citeproctext]{ref-Persico2000}
Persico, Nicola. 2000. {``Information Acquisition in Auctions.''}
\emph{Econometrica} 68 (1): 135--48.
\url{https://doi.org/10.1111/1468-0262.00096}.

\bibitem[\citeproctext]{ref-RobertsSweeting2013}
Roberts, James W., and Andrew Sweeting. 2013. {``When Should Sellers Use
Auctions?''} \emph{American Economic Review} 103 (5): 1830--61.
\url{https://doi.org/10.1257/aer.103.5.1830}.

\end{CSLReferences}

\section{Paper appendix}\label{paper-appendix}

\subsection{A.1. Acquisition payoffs and Proposition
1}\label{pa-payoffs}

\begin{enumerate}
\def\labelenumi{\arabic{enumi}.}
\tightlist
\item
  \textbf{Determine payments before taking expectations.} Without entry,
  payment is \(p\mathbf1\{R\ge p\}\). With high quality, the challenger
  wins and pays \(\max\{p,R\}\). With low quality, target proceeds are
  \(\max\{p,\min(R,\ell)\}\) and the challenger obtains
  \((\ell-\max\{p,R\})_+\). Truthful bidding is weakly dominant
  conditional on every competing bid, so these payoffs do not depend on
  a conjectured shading strategy.
\item
  \textbf{Integrate the uniform realization.} We obtain
\end{enumerate}

\[
\begin{aligned}
t_H&=\frac1r\left[\int_0^p p\,du+\int_p^r u\,du\right],\\
t_L&=\frac1r\left[\int_0^p p\,du+\int_p^\ell u\,du+\int_\ell^r\ell\,du\right],\\
g_L&=\frac1r\left[\int_0^p(\ell-p)\,du+\int_p^\ell(\ell-u)\,du\right].
\end{aligned}
\tag{A.1}
\]

Evaluating gives (4), and differentiating gives (5). 3.
\textbf{Establish the distribution-free opposition.} If \(R\le\ell\),
high and low target payments coincide. If \(R>\ell\), their difference
is \(R-\ell\). Hence the difference is \((R-\ell)_+\). Tonelli's theorem
applied to nonnegative indicators gives both integrals in (6). A
stronger distribution has a smaller CDF, increasing the survival
integral and decreasing the profit integral. Positive integral
differences give strictness. A fixed-posterior weighted average
preserves the profit ordering. Online Appendix A.2 provides the measure
formulation.

\subsection{A.2. Theorem 1: inference and global equilibrium
verification}\label{pa-entry}

\begin{enumerate}
\def\labelenumi{\arabic{enumi}.}
\tightlist
\item
  \textbf{Bound beliefs under all mixed orders.} Positivity of \(f\)
  makes every conditional order-flow density positive. Integrating the
  pairwise inequality (8) gives \(m\le\mu_X\le M\). Conditional
  expectation gives the same bound for beliefs based on price. Both
  \(g_H\) and \(g_L\) decrease with \(r\), so (A1) implies
  \(c_L<B_r(m)\) in both economies.
\item
  \textbf{Recover the buyer's actual information.} The low-cost type
  enters at every feasible price, so \(e(P)\ge\rho\). Conditional
  pricing gives (11), whose positive denominator makes \(\mu_X\)
  measurable with respect to \(P\). Therefore \(\Pr(H\mid P)=\mu_X\)
  almost surely. The conditional-payoff subtraction yields (10). This
  uses independence of \(R,C\) from trading, not an assumption that the
  buyer sees order flow.
\item
  \textbf{Eliminate every nonzero weak-economy order.} A correctly
  signed order with magnitude \(s>0\) yields at most
  \(s[\Delta_T(r_0)-k]<0\). A wrong-signed order has negative gross
  payoff and also pays its cost. Thus zero is the only best response
  against every candidate equilibrium, including mixed ones. Under zero
  orders, \(\mu_X=1/2\), and (A1)--(A2) give entry \(\rho\). Constant
  pricing constructs that equilibrium.
\item
  \textbf{Control the entire strong-economy order interval.} For a
  bounded residual \(A\), the translation formula
\end{enumerate}

\[
F(s_2)-F(s_1)
=\int_{s_1}^{s_2}\int \partial_u f(x\mp u)A(x)\,dx\,du
\tag{A.2}
\]

follows by the fundamental theorem for the absolutely continuous density
and Fubini: the absolute double integral is bounded by
\(\|A\|_\infty|s_2-s_1|\|f'\|_1\). Thus \(F\) is absolutely continuous,
with \(|F'|\le F/b\) almost everywhere. Equation (10) gives
\(F\ge\rho m\Delta_T\). Differentiating \(U=sF-ks\) and applying (A3)
yields (13). Integrating the positive lower derivative bound shows that
the full correctly signed order strictly dominates every smaller
magnitude. Wrong signs are dominated by zero. Hence every equilibrium
has full correctly signed orders; mixing cannot introduce another
optimal action. 5. \textbf{Construct and verify the informative
equilibrium.} Full orders give

\[
\mu_X(x)=\left[1+\exp\left\{-\frac{|x+1|-|x-1|}{b}\right\}\right]^{-1}.
\tag{A.3}
\]

Define entry and price by (9). The strict monotonicity of \(P_r(\mu)\)
gives a measurable inverse on its image, including across its upward
jump. The buyer therefore recovers the required posterior from price,
optimizes preparation, and subsequently bids truthfully. The preceding
global bounds verify investor optimality and market-maker pricing,
proving existence and the stated uniqueness. 6. \textbf{Compute entry,
ownership, and the information comparison.} Since \(B_r\) increases with
\(\mu\) and decreases with \(r\), (A2) gives \(1/2<\tau<M\). In the
central likelihood region, solving \(\mu_X=\tau\) gives \(x^*\) in (14).
The Laplace survival function at \(x^*-1\) and \(x^*+1\) gives
\(\alpha_H,\alpha_L\) and the entry and ownership formulas. The
favorable event has positive probability. Ignoring the strong-economy
price reproduces the uninformative experiment; no state-independent
transformation of a constant experiment can reproduce its nonconstant
state-dependent law. The strong experiment therefore strictly Blackwell
dominates the weak one. Strict primitive margins establish an open
region. Online Appendix A.1--A.4 supplies conditional-probability
versions, null-set invariance under deviations, and full convolution
regularity.

\subsection{A.3. Fixed information, thresholds, and Corollary
3}\label{pa-thresholds}

\begin{enumerate}
\def\labelenumi{\arabic{enumi}.}
\tightlist
\item
  \textbf{Fixed experiment.} Couple the economies using the same
  realization of the signal and cost. Since
  \(B_{r_1}(\mu)\le B_{r_0}(\mu)\), the entry indicator is pointwise
  weakly smaller at the stronger rival. A strict decrease in expected
  entry occurs exactly when the set
  \(\{B_{r_1}(\mu)<C\le B_{r_0}(\mu)\}\) has positive probability. This
  proves Proposition R1 under the stated entry-at-indifference
  convention.
\item
  \textbf{Exact pooling existence.} In the domain of Proposition 3,
  pooling implies entry \(\rho\) and constant price. A correctly signed
  deviating order earns
\end{enumerate}

\[
s\left(\frac{\rho\Delta_T(r)}2-k\right).
\tag{A.4}
\]

Thus zero is optimal exactly when its coefficient is nonpositive.
Solving \(\Delta_T(r)=d\) on \(r>\ell\) gives \(\mathfrak r(d)\). This
proves the exact \(r_N\) boundary. In contrast, \(\Delta_T<k\)
eliminates nonzero orders against every candidate schedule and
establishes the smaller sufficient uniqueness region. 3. \textbf{Full
orders and entry feasibility.} The derivative bound (13) is uniform
across candidate equilibria whenever \(r>r_U\) and the low-cost floor
holds. For expensive entry, the largest feasible belief is \(M\). Its
profitability is

\[
B_r(M)=Mh-\frac{Mr}{2}+\frac{(1-M)\ell^2-p^2}{2r}.
\tag{A.5}
\]

It is strictly decreasing on the auction-support domain. Multiplying
\(B_r(M)=c_H\) by \(2r\) gives a quadratic; selecting a root within that
domain gives (16). When the discriminant or domain requirement fails,
the feasibility comparison is made directly rather than assigning an
artificial crossing. The formula in (16) selects the larger algebraic
root; any other root cannot provide another crossing in the specified
domain because of strict monotonicity. 4. \textbf{Endpoint behavior.}
Under full Laplace orders, \(\mu_X=M\) for \(x\ge1\), an event with
positive probability. Our tie rule admits expensive entry at \(r_C\) on
that event. As \(r\) approaches \(r_C\) from below,

\[
\mathsf E(r)\longrightarrow\rho+\frac{1-\rho}{4}(1+e^{-2/b}).
\tag{A.6}
\]

For strengths strictly above \(r_C\), expensive entry is impossible.
With logistic noise the posterior reaches \(M\) only at an infinite flow
limit, so expensive entry instead converges to zero continuously. These
statements concern their respective information experiments. 5.
\textbf{The finite nonmonotonicity corollary.} Apply Theorem 1 at
\(r_0,r_1\). At \(r_2\), the low-cost floor and full-order derivative
bound establish unique full orders, while \(B_{r_2}(M)<c_H\) excludes
every expensive entrant. Hence entry returns to \(\rho\). This proves
(17) without assigning a unique branch at intermediate strengths. Full
derivations are in Online Appendix A.5.

\subsection{A.4. The computer-assisted asymmetric
result}\label{pa-certificate}

\begin{enumerate}
\def\labelenumi{\arabic{enumi}.}
\tightlist
\item
  \textbf{Candidate beliefs and entry.} Fix \((q_H,q_L)=(1,-v)\),
  \(0<v<1\). Bayes' rule gives
  \(\mu_v(x)=\operatorname{logistic}((|x+v|-|x-1|)/b)\). The posterior
  ranges from \(m_v=(1+e^{(1+v)/b})^{-1}\) to \(M_v=1-m_v\). In the
  certified region \(1/2<\tau<M_v\), and
\end{enumerate}

\[
x^*(r,v)=\frac{b\operatorname{logit}(\tau)+1-v}{2},\qquad
\mathsf E(r,v)=\rho+\frac{1-\rho}{2}
\left[1-\frac12e^{(x^*-1)/b}+\frac12e^{-(x^*+v)/b}\right].
\tag{A.7}
\]

\begin{enumerate}
\def\labelenumi{\arabic{enumi}.}
\setcounter{enumi}{1}
\tightlist
\item
  \textbf{A globally optimal interior short.} The residual
  \(A_L=e(\mu_v)\Delta_T\mu_v\) is bounded, nonconstant, and
  nondecreasing. Its finite positive Stieltjes measure gives
\end{enumerate}

\[
F_L'(s)=-\int f(x+s)\,dA_L(x)<0,
\qquad |F_L''(s)|\le-\frac1bF_L'(s).
\tag{A.8}
\]

Therefore \(U_L''(s)\le(2-s/b)F_L'(s)<0\) for \(s\in[0,1]\) when
\(b>1/2\). The measure includes the entry jump; no derivative of that
jump is omitted. A root of the unilateral marginal-profit equation is
thus the unique global short magnitude against its candidate schedule.
3. \textbf{An exact root, not a small residual.} Let
\(\Psi(r,v)=U_L'(v;1,-v)\), with the unilateral derivative taken holding
the candidate schedule fixed. Recomputing the schedule as the candidate
parameter \(v\) varies gives a continuous \(\Psi\) on each certified
bracket. Outward interval evaluation proves positive \(\Psi\) at the
left endpoint and negative \(\Psi\) at the right endpoint. The
intermediate value theorem gives an exact root inside. Throughout each
bracket, the entry threshold stays strictly between \(-v\) and \(1\).
The endpoints \(v_{j,-},v_{j,+}\) are those of the corresponding bracket
in (18). The enclosures completing the sign tests are

\[
\begin{aligned}
\Psi(r_{\mathrm a},v_{\mathrm a,-})&\ge0.000000000114>0,&
\Psi(r_{\mathrm a},v_{\mathrm a,+})&\le-0.000000000096<0,\\
\Psi(r_{\mathrm b},v_{\mathrm b,-})&\ge0.000000000123>0,&
\Psi(r_{\mathrm b},v_{\mathrm b,+})&\le-0.000000000121<0,\\
\Psi(r_{\mathrm c},v_{\mathrm c,-})&\ge0.000000000172>0,&
\Psi(r_{\mathrm c},v_{\mathrm c,+})&\le-0.000000000242<0.
\end{aligned}
\tag{A.9}
\]

\begin{enumerate}
\def\labelenumi{\arabic{enumi}.}
\setcounter{enumi}{3}
\tightlist
\item
  \textbf{Exclude every smaller purchase.} For any bounded residual
  \(A_H\in[0,\Delta_T]\), the Laplace kernel satisfies
  \(f''=(f-\delta_0)/b^2\) in distributions. Consequently
\end{enumerate}

\[
F_H''(s)=\frac{F_H(s)-A_H(s)}{b^2}\quad\text{a.e.},\qquad
|U_H''(s)|\le L_U:=\frac{2\Delta_T}{b}+\frac{\Delta_T}{b^2}.
\tag{A.10}
\]

On the mesh \(s_j=j/n\), interval arithmetic bounds \(U_H'(s_j)\)
uniformly over the entire root bracket. Every untested magnitude is
within \(1/(2n)\) of a mesh point, so

\[
\inf_{s\in[0,1]}U_H'(s)
\ge\min_j\underline{U_H'(s_j)}-\frac{L_U}{2n}>0.
\tag{A.11}
\]

The certified global lower margins at the ordered strengths are
0.0000761777, 0.0027531948, and 0.0054921767. Thus the favorable type
chooses its maximum purchase. All wrong-signed trades have negative
gross profit and are dominated by zero. 5. \textbf{Integrate exactly and
conclude.} Split each convolution at \(-v,x^*,1\), and the deviation
center. On the central region put \(c=(1-v)/2\) and \(t=e^{(x-c)/b}\),
so \(\mu_v=t^2/(1+t^2)\). The primitives of \(e^{x/b}(1-\mu_v)\),
\(e^{-x/b}(1-\mu_v)\), \(e^{x/b}\mu_v\), and \(e^{-x/b}\mu_v\) are
respectively

\[
b e^{c/b}\arctan t,\quad
b e^{-c/b}(-t^{-1}-\arctan t),\quad
b e^{c/b}(t-\arctan t),\quad
b e^{-c/b}\arctan t.
\tag{A.12}
\]

Constant-posterior tails integrate as exponentials. Interval arithmetic
applied to these expressions encloses the root tests, (A.11), and the
entry probabilities. The pooling margins are positive at the same
strengths, and the entry enclosures in (18) are disjoint and ordered.
This proves Proposition 4. Online Appendix B specifies the interval
arithmetic, parameter brackets, endpoint ordering, and every regularity
argument needed for replication.

\subsection{A.5. Smooth distributions and moderate acquisition
values}\label{pa-smooth}

\begin{enumerate}
\def\labelenumi{\arabic{enumi}.}
\tightlist
\item
  \textbf{Log-density control.} If a positive density is absolutely
  continuous and \(|f'|\le Lf\) almost everywhere with \(L<1\),
  log-density integration gives the pairwise likelihood bound
  \(e^{-2L}\le f(x-q)/f(x-q')\le e^{2L}\). The convolution proof in A.2
  yields \(U'(s)\ge(1-L)\rho m_L\Delta_T-k\), where
  \(m_L=(1+e^{2L})^{-1}\). Logistic noise has
  \((\log f)'=-\tanh(z/(2b))/b\) and therefore \(L=1/b\).
\item
  \textbf{Attainable informative beliefs.} Under full logistic orders
  the posterior log odds are
\end{enumerate}

\[
2\log\cosh\left(\frac{x+1}{2b}\right)
-2\log\cosh\left(\frac{x-1}{2b}\right).
\tag{A.13}
\]

Its derivative is positive and its limits are \(-2/b,2/b\). It spans the
posterior interval \((m,M)\). Solving for the interior threshold gives
(20); the positive density assigns positive probability above it. The
uniform no-trade and full-order arguments and the ownership/Blackwell
comparisons follow as in Theorem 1. This proves Corollary 1. 3.
\textbf{Atomless preparation costs.} Under (21), all low-cost
realizations participate at all feasible prices, preserving the residual
lower bound. At the weak prior all high-cost realizations stay out. At
prices sufficiently close to the upper feasible posterior all high-cost
realizations enter. That event has positive probability under either
noise law. The cost CDF is nondecreasing, so the price construction
remains strictly increasing. The global trading bounds are unchanged.
This proves Corollary 2. 4. \textbf{Nonemptiness with any strict value
gap.} At \(r=\ell\), \(g_H-g_L=h-\ell\) and

\[
B_\ell(M)-B_\ell(1/2)=(M-1/2)(h-\ell)>0.
\tag{A.14}
\]

Choose \(r_1>\ell\) sufficiently close that \(B_{r_1}(M)>B_\ell(1/2)\).
Next choose \(r_0\in(\ell,r_1)\) sufficiently close to \(\ell\) that
\(\Delta_T(r_0)<(1-1/b)\rho m\Delta_T(r_1)\). Choose \(k\) strictly
between these bounds, \(c_H\) strictly between \(B_{r_0}(1/2)\) and
\(B_{r_1}(M)\), and \(c_L\) positive and below \(B_{r_1}(m)\). These
intervals are nonempty. Continuity preserves all strict inequalities in
a neighborhood. Online Appendix A.6 supplies full regularity, the
logistic inversion, and further prior/cost extensions.

\subsection{A.6. Theorem R2: the buyer and trader have distinct private
signals}\label{pa-signals}

\begin{enumerate}
\def\labelenumi{\arabic{enumi}.}
\tightlist
\item
  \textbf{Control public and private posteriors.} Marginal trader
  signals are equally likely. Arbitrary orders mixed conditional on
  those signals obey the same likelihood-ratio bound as in Lemma 1, so
  \(m\le\lambda_X\le M\). Conditional independence gives
  \(\mu_X=(1-a)+(2a-1)\lambda_X\) and the public interval
  \([\mu_-,\mu_+]\). The buyer's posterior is \(\phi_Y(\mu_P)\) because
  \(Y\) is independent of \(X\) conditional on fundamental quality. Its
  lowest feasible posterior is \(\phi_-(\mu_-)\). The first restriction
  in (27) makes low-cost entry optimal in both economies.
\item
  \textbf{Price state-dependent participation.} Since
  \(\phi_+\ge\phi_-\), optimal high-cost entry satisfies \(I_+\ge I_-\).
  Formula (24) follows by averaging over \(Y\) conditional on \(H\) or
  \(L\). Thus \(e_H\ge e_L\ge\rho\), and
\end{enumerate}

\[
D=e_L\Delta_T+(e_H-e_L)w_H,
\quad 0\le e_H-e_L\le(1-\rho)(2d-1).
\tag{A.15}
\]

This proves the bounds in (24). Equation (25) then makes the public
posterior measurable with respect to price. For construction, entry is
nondecreasing in the public posterior for each private signal. Between
entry changes the price slope is \(D>0\); at an entry change the price
increases because the added conditional target payoffs \(w_H,w_L\) are
positive. The buyer can therefore recover public information and combine
it with \(Y\). 3. \textbf{Derive residuals for the trader's
information.} Conditional on \(T=+\), the probability of \(H\) is \(a\)
even after observing any flow generated by a signal-contingent
unilateral order; for \(T=-\) it is \(1-a\). This follows because order
randomization and noise add no information about \(\theta\) conditional
on \(T\). Subtracting the competitive price from the appropriate
conditional target payoff gives (26), with

\[
m(2a-1)\rho\Delta_T\le A_+,A_-
\le(2a-1)[\Delta_T+(1-\rho)(2d-1)w_H].
\tag{A.16}
\]

\begin{enumerate}
\def\labelenumi{\arabic{enumi}.}
\setcounter{enumi}{3}
\tightlist
\item
  \textbf{Establish necessity and existence in both economies.} The
  lower cost inequality in (27) makes the upper residual at \(r_0\)
  smaller than \(k\), excluding every nonzero correctly signed order;
  wrong signs are dominated. Public beliefs are then the prior, and the
  buyer's highest private posterior is \(d\). Since \(B_{r_0}(d)<c_H\),
  entry is \(\rho\) and constant pricing constructs the equilibrium. At
  \(r_1\), the convolution argument gives
\end{enumerate}

\[
U'(s)\ge(1-1/b)m(2a-1)\rho\Delta_T(r_1)-k>0.
\tag{A.17}
\]

Full correctly signed orders are necessary against every candidate
schedule. Bayes' rule, optimal private-signal entry, and the strictly
increasing price construction establish existence. Finally,
\(c_H<B_{r_1}(\phi_+(\mu_+))\) makes expensive entry worthwhile for a
favorable private signal and a positive-probability set of sufficiently
favorable prices. This proves the strict entry comparison. Online
Appendix A.7 gives the joint conditional laws, measurable construction,
and complete probability formulas.

\subsection{A.7. Bargaining and Lemma 3}\label{pa-bargaining}

\begin{enumerate}
\def\labelenumi{\arabic{enumi}.}
\tightlist
\item
  \textbf{Derive the price from a feasible fallback.} Let \(V\) be the
  winner's value and \(z\) the next-best buyer's value. The seller's
  disagreement payoff is \(z\) because the fallback sale is enforceable.
  The winner's disagreement payoff is zero. The Nash solution maximizes
  \((P-z)^\eta(V-P)^{1-\eta}\) on \([z,V]\). For interior weights its
  first-order condition gives \(P=z+\eta(V-z)\); continuity covers the
  zero-weight endpoint.
\item
  \textbf{Allocate the challenger surplus.} A challenger wins only if
  \(\theta\ge R\) and then keeps \((1-\eta)(\theta-R)\). Hence
  \(G_{\theta,\eta}\) in (28). Without entry the same institution has
  \(z=0\), giving target payment \(\eta R\).
\item
  \textbf{Subtract the target payments state by state.} If \(R\le\ell\),
  the high--low difference is \(\eta(h-\ell)\). If \(R>\ell\), it is
  \(\eta(h-\ell)+(1-2\eta)(R-\ell)\). Taking expectations proves (28).
  FOSD applies to the increasing function \((R-\ell)_+\) and the
  decreasing function \((\theta-R)_+\). The coefficient \(1-2\eta\)
  determines the information-spread ordering. These are payment-stage
  conclusions for the specified verifiable-value institution. Online
  Appendix A.8 supplies the full bargaining construction and endpoint
  treatment.
\end{enumerate}

\subsection{A.8. Proposition 2: a matched surplus
comparison}\label{pa-welfare}

\begin{enumerate}
\def\labelenumi{\arabic{enumi}.}
\tightlist
\item
  \textbf{Resolve the price-hidden game.} The buyer's posterior stays at
  the prior, so (A1)--(A2) give entry \(\rho\). The strong global bound
  still applies with constant \(e=\rho\), establishing full orders. The
  same investor order magnitudes and noise distribution therefore occur
  with and without feedback.
\item
  \textbf{Compute incremental allocation value pointwise.} If \(R<p\),
  entry replaces no sale by a sale to \(\theta\), and the right-hand
  side of (29) equals \(\theta-p+p=\theta\). If \(R\ge p\), it equals
  \((\theta-R)_+\), the improvement over incumbent ownership. This
  proves (29). Integrating over \(R\) gives
  \(g_\theta+p\Pr(R<p)=g_\theta+p^2/r\).
\item
  \textbf{Condition on the information used to enter.} Every additional
  entrant satisfies \(B_r(\mu)-C\ge0\). Its conditional net contribution
  is therefore at least \(p^2/r>0\). The additional-entry event has
  positive probability; low-cost participation is unchanged. Averaging
  gives the first expression in (30), or its integral over high costs
  for Corollary 2. Noise-trader transfers and acquisition payments do
  not alter total allocation surplus; identical real trading costs
  cancel.
\item
  \textbf{Compare target proceeds and match price levels.} Since
  \(t_H,t_L>t_0\), additional entry raises target proceeds by the second
  expression in (30). Adding a deterministic external dividend to the
  price-hidden traded claim increases its rational price by the same
  amount and leaves \(V_T-P\) unchanged. It therefore preserves trading
  and entry. Online Appendix A.9 states the probability coupling and the
  cost-distribution integral explicitly.
\end{enumerate}

\subsection{A.9. Numerical primitive vectors}\label{pa-parameters}

The benchmark parameter vector is

\[
\begin{aligned}
(h,\ell,p,\rho,c_L,c_H,b,k)
&=(10,1,0.5,0.25,1,6,2,0.02),\\
(r_0,r_1,r_2)&=(1.2,3,3.6).
\end{aligned}
\tag{A.18}
\]

The moderate-value vector is

\[
\begin{aligned}
(h,\ell,p,\rho,c_L,c_H,b,k)
&=(2,1,0.5,0.25,0.3,0.89,2,0.002),\\
(r_0,r_1)&=(1.05,1.5).
\end{aligned}
\tag{A.19}
\]

For complementary signals we use

\[
\begin{aligned}
(h,\ell,p,\rho,c_L,c_H,b,k)
&=(10,1,0.5,0.85,1,7.14,2,0.015),\\
(r_0,r_1,a,d)
&=(1.1,2.3,0.70,0.75).
\end{aligned}
\tag{A.20}
\]

The atomless-cost half-width is \(\varepsilon_C=0.1\) and the
atomless-value half-width is \(\varepsilon_V=0.05\). These are separate
experiments. Online Appendix C defines exact input values, derived
quantities, formatting, and acceptance criteria for every placeholder.

\subsection{A.10. Sale-domain identities and local
differentiation}\label{pa-design}

\begin{enumerate}
\def\labelenumi{\arabic{enumi}.}
\tightlist
\item
  \textbf{Condition on reserve admissibility.} When \(v<p\), the
  challenger never meets the reserve, so its presence leaves target
  proceeds at \(t_0\) and its acquisition profit at zero. When
  \(v\ge p\), the target sells to the larger of \(R,v\) and receives
  \(\max\{p,\min(R,v)\}\). These cases prove (33), including the stated
  equality convention. Averaging over \(v\) gives conditional class
  payoffs for atomless values.
\item
  \textbf{Derive seller revenue.} Conditional on challenger quality,
  entry changes target proceeds by \(t_\theta-t_0\). Independence of
  \(R\) from entry justifies (31). Within \(0<p<\ell\),
  \(\partial_pt_0=1-2p/r\), \(\partial_pt_H=\partial_pt_L=p/r\), and
  \(\partial_pB_{p,r}(\mu)=-p/r\) at a fixed posterior. Applying the
  product rule to (31) gives (34).
\item
  \textbf{Differentiate entry only on a regular branch.} Assume an
  atomless cost CDF that is continuously differentiable on the reached
  profit range, differentiable state-contingent orders and posteriors
  almost everywhere in the noise realization, and an integrable common
  bound for the differentiated entry integrand in a neighborhood of the
  reserve. Dominated differentiation of
  \(e_\theta=\int f(z)H_C(B_{p,r}(\mu_\theta(z;p)))dz\) yields (35). At
  reserve/value boundaries, nondifferentiable order branches, or atomic
  cost thresholds, use the level objective rather than this derivative.
\item
  \textbf{Keep the optimization problem separate.} A valid seller
  optimum must compare continuations after all feasible reserve choices.
  The payment identities and local decomposition do not select a
  continuation or prove attainment of that optimum. Online Appendix A.10
  gives the continuous-class formulas and sufficient differentiation
  conditions; Appendix C.6 specifies the exploratory correspondence
  calculation that supports the next theorem.
\end{enumerate}

\end{document}
